\documentclass[../../main.tex]{subfiles}

\begin{document}

\chapter{Layanan Keamanan dan Serangan Kriptografi}

\begin{learningoutcome}
Setelah mempelajari bab ini, mahasiswa diharapkan mampu:
\begin{itemize}
    \item Menjelaskan konsep dasar serangan terhadap sistem kriptografi dan tujuan dari serangan tersebut (Sub-CPMK 2.1).
    \item Mengidentifikasi dan membedakan berbagai jenis serangan kriptografi berdasarkan informasi yang tersedia (Ciphertext-only, Known-plaintext, Chosen-plaintext, dll) (Sub-CPMK 2.1).
    \item Membedakan antara serangan pasif dan serangan aktif serta memberikan contoh masing-masing (Sub-CPMK 2.1).
    \item Memahami konsep keamanan komputasi (\textit{computationally secure}) dan serangan side-channel (Sub-CPMK 2.1).
\end{itemize}
\end{learningoutcome}

\section{Konsep Dasar Serangan}
Tujuan utama dari kriptografi adalah menjaga kerahasiaan pesan atau kunci dari pihak ketiga, yang dalam konteks ini disebut sebagai musuh, lawan, penyadap, atau kriptanalis.

Sebuah \textbf{serangan (\textit{attack})} diartikan sebagai setiap usaha atau percobaan yang dilakukan oleh pihak lawan untuk menemukan kunci atau menemukan plainteks dari cipherteksnya. Dalam menganalisis serangan, kita menggunakan \textbf{Prinsip Kerckhoffs}:
\begin{quote}
    \textit{Semua algoritma kriptografi harus publik; hanya kunci yang rahasia.}
\end{quote}
Oleh karena itu, keamanan sistem kriptografi tidak terletak pada kerahasiaan algoritma, melainkan sepenuhnya terletak pada kerahasiaan \textbf{kunci}.

\section{Jenis-Jenis Serangan Berdasarkan Informasi}
Berdasarkan informasi yang dimiliki oleh penyerang, serangan dapat dikategorikan menjadi beberapa jenis:

\subsection{Ciphertext-Only Attack}
Ini adalah serangan yang paling sulit karena kriptanalis hanya memiliki akses ke cipherteks saja. Teknik yang digunakan biasanya meliputi:
\begin{itemize}
    \item \textit{Exhaustive key search} (Brute force).
    \item Analisis frekuensi karakter.
\end{itemize}
Contoh: Pada Caesar Cipher, jika penyerang hanya memiliki cipherteks \texttt{KHOORZRUOG}, ia dapat mencoba semua 25 kemungkinan pergeseran huruf hingga menemukan plainteks yang masuk akal (\texttt{HELLOWORLD}).

\subsection{Known-Plaintext Attack}
Penyerang memiliki sejumlah pasangan plainteks dan cipherteks yang berkoresponden ($P_i$ dan $C_i = E_k(P_i)$). Serangan ini sering terjadi jika sebagian isi pesan dapat ditebak, seperti \textit{header} file, protokol standar, atau salam pembuka email (misalnya "Dengan hormat").

\subsection{Chosen-Plaintext Attack}
Penyerang dapat memilih plainteks tertentu untuk dienkripsi oleh sistem dan mendapatkan cipherteksnya. Dengan memilih plainteks yang strategis (misalnya rangkaian huruf 'AAAAA'), penyerang dapat lebih mudah mendeduksi kunci.

\subsection{Adaptive Chosen-Plaintext Attack}
Variasi dari chosen-plaintext attack di mana penyerang dapat menyesuaikan plainteks berikutnya berdasarkan hasil cipherteks dari query sebelumnya. Penyerang belajar secara bertahap untuk menggali informasi kunci.

\subsection{Chosen-Ciphertext Attack}
Penyerang dapat memilih cipherteks tertentu dan mendapatkan hasil dekripsinya melalui sebuah \textit{oracle} (sistem yang bisa mendekripsi). Informasi ini kemudian digunakan untuk memecahkan cipher atau menemukan kunci.

\section{Tujuan Serangan (Hasil yang Ingin Dicapai)}
Selain berdasarkan informasi, serangan juga dapat dibedakan berdasarkan target akhirnya:
\begin{itemize}
    \item \textbf{Key-Recovery Attack}: Tujuan utamanya adalah memperoleh kunci kriptografi.
    \item \textbf{Plaintext-Recovery Attack}: Tujuan utamanya adalah memperoleh plainteks, meskipun kunci tidak ditemukan.
    \item \textbf{Distinguishing Attack}: Menentukan apakah suatu data adalah keluaran algoritma kriptografi atau data acak (\textit{random}).
    \item \textbf{Forgery Attack}: Membuat pesan palsu yang terlihat sah bagi sistem (sangat krusial pada MAC dan tanda tangan digital).
    \item \textbf{Replay Attack}: Mengirim kembali pesan sah yang pernah dicegat sebelumnya untuk memicu tindakan berulang tanpa perlu memecahkan enkripsi.
\end{itemize}

\section{Serangan Pasif vs Serangan Aktif}
Berdasarkan keterlibatan penyerang dalam komunikasi:

\subsection{Serangan Pasif (\textit{Passive Attack})}
Penyerang tidak mengubah data atau terlibat aktif dalam komunikasi; ia hanya menyadap informasi.
\begin{itemize}
    \item \textbf{Interception}: Menyadap konten pesan.
    \item \textbf{Traffic Analysis}: Mengamati pola lalu lintas data (siapa berkomunikasi dengan siapa, kapan, dan berapa banyak data yang dikirim).
\end{itemize}
Contoh alat: \textbf{Wireshark} dapat digunakan untuk memantau lalu lintas jaringan dan menemukan plainteks jika data tidak dienkripsi.

\subsection{Serangan Aktif (\textit{Active Attack})}
Penyerang mengintervensi komunikasi dan memodifikasi sistem.
\begin{itemize}
    \item \textbf{Interruption}: Memutus koneksi sehingga layanan tidak tersedia (DoS).
    \item \textbf{Fabrication}: Mengirimkan pesan palsu yang seolah-olah sah.
    \item \textbf{Modification}: Mengubah isi pesan di tengah jalan.
    \item \textbf{Replay}: Mengirim ulang pesan sah.
    \item \textbf{Man-in-the-Middle (MITM)}: Penyerang menempatkan dirinya di antara pengirim dan penerima, mengintersepsi dan memodifikasi komunikasi kedua belah pihak secara real-time.
\end{itemize}

\section{Keamanan Komputasi dan Side-Channel Attack}
Sebuah algoritma dikatakan \textbf{aman secara komputasi (\textit{computationally secure})} jika:
1. Persamaan matematikanya terlalu kompleks untuk dipecahkan secara analitik.
2. Biaya untuk memecahkan cipherteks melebihi nilai informasi di dalamnya.
3. Waktu yang dibutuhkan untuk memecahkan cipherteks melebihi masa berlaku kerahasiaan informasi tersebut.

\subsection{Side-Channel Attack}
Serangan ini tidak menyerang algoritma matematika, melainkan memanfaatkan kebocoran fisik saat algoritma dijalankan pada perangkat:
\begin{itemize}
    \item \textbf{Timing Attack}: Menganalisis perbedaan waktu eksekusi.
    \item \textbf{Power Analysis}: Menganalisis konsumsi daya (SPA/DPA).
    \item \textbf{Electromagnetic Attack}: Menganalisis radiasi elektromagnetik.
    \item \textbf{Fault Injection}: Sengaja membuat kesalahan pada perangkat untuk menganalisis output yang salah (DFA).
\end{itemize}

\begin{obeactivity}
\textbf{Aktivitas 2.1: Simulasi Analisis Paket dengan Wireshark}
Gunakan aplikasi Wireshark untuk menangkap paket data dari protokol yang tidak terenkripsi (misalnya HTTP biasa). Cobalah cari paket \textit{POST} dan lihat apakah Anda dapat menemukan \textit{username} atau \textit{password} dalam bentuk plainteks. Diskusikan mengapa penggunaan HTTPS sangat krusial untuk mencegah serangan pasif.
\end{obeactivity}

\begin{obereflection}
\textbf{Latihan 2.1}
1. Mengapa \textit{Chosen-Plaintext Attack} dianggap lebih berbahaya daripada \textit{Ciphertext-Only Attack}? Berikan contoh skenario di mana penyerang bisa melakukan serangan ini.
2. Berikan contoh nyata dari \textit{Replay Attack} dalam sistem perbankan online dan jelaskan bagaimana penggunaan \textit{nonce} atau \textit{timestamp} dapat mencegahnya.
\end{obereflection}

\begin{obeassessment}
\textbf{Kuis Bab 2}
1. Apa perbedaan mendasar antara serangan pasif dan serangan aktif? Manakah yang lebih sulit dideteksi? Jelaskan alasannya.
2. Dalam kondisi apa sebuah algoritma yang secara matematis kuat tetap bisa ditembus menggunakan \textit{Side-Channel Attack}?
\end{obeassessment}

\begin{competencychecklist}
\begin{tabularx}{\linewidth}{|X|c|}
\hline
Indikator Kompetensi & Tercapai \\
\hline
Saya dapat menjelaskan prinsip Kerckhoffs dan mengapa kunci harus rahasia & [ ] \\
\hline
Saya dapat membedakan jenis serangan berdasarkan informasi (COA, KPA, CPA, CCA) & [ ] \\
\hline
Saya dapat mengidentifikasi perbedaan antara serangan pasif dan aktif & [ ] \\
\hline
Saya dapat menjelaskan konsep keamanan komputasi dan side-channel attack & [ ] \\
\hline
\end{tabularx}
\end{competencychecklist}

\end{document}
